\documentclass[../main.tex]{subfiles}

\begin{document}
	
	\chapter{Metodi e dataset}
	\label{ch:metodi_dataset}
	
	\section{Il dataset originale}
	\label{sec:dataset_originale}
	I dati viscosimetrici impiegati in questo studio derivano dal database pubblico SciGlass (\url{https://github.com/epam/SciGlass}) e si focalizzano esclusivamente sui liquidi formatori di vetro a base di ossidi. A causa della difficoltà incontrata nel reperire direttamente i dati, si è assunto come punto di partenza il dataset già estratto, deduplicato e pulito da Cassar~\cite{cassarViscNetNeuralNetwork2021}. Tale procedura preliminare prevedeva l'esclusione di misurazioni estreme (viscosità superiori a $10^{12}$ Pa$\cdot$s o inferiori a $10^{-5}$ Pa$\cdot$s) e la rimozione di liquidi privi di un numero sufficiente di misurazioni nelle regioni di alta e bassa viscosità.
	
	Tuttavia, l'approccio di deduplicazione adottato originariamente presenta una criticità metodologica non trascurabile. La procedura raggruppa i campioni aventi la medesima composizione chimica nominale sovrapponendo e raggruppando dati ottenuti da fonti sperimentali diverse, le quali possono aver utilizzato protocolli di fusione, velocità di raffreddamento o tecniche di misurazione disomogenee. Di conseguenza, per alcuni liquidi si assiste a una dispersione artificiale dei dati che genera difficoltà in fase di \textit{fitting}, introducendo una varianza che non deriva dalla termodinamica del sistema, bensì dall'incompatibilità intrinseca dei dati sperimentali storicamente accorpati (vedi Fig.~\ref{fig:deduplication}).
	
	\subsection*{Suddivisione del dataset}
	\label{sec:data_split}
	In seguito alla pulizia e deduplicazione dei dati di SciGlass fatta da Cassar, il dataset finale utilizzabile risulta composto da 847 differenti liquidi a base di ossido, per un totale di 17584 punti dati viscosimetrici misurati~\cite{cassarViscNetNeuralNetwork2021}. 
	
	La strategia di partizionamento dei dati riveste un ruolo importante per la validazione dei modelli. Per questo lavoro si è scelto di mantenere l'esatto \textit{split} dei dati definito in precedenza per ViscNet permettendo di svolgere un confronto diretto con i risultati ottenuti dai modelli lineari analizzati. Degli 847 liquidi disponibili, quindi, un sottoinsieme di 85 liquidi è stato selezionato casualmente da Cassar per formare l'\textit{unseen test dataset}. Le misurazioni appartenenti a questi 85 liquidi non sono mai state esposte ai modelli durante le fasi di ottimizzazione degli iperparametri o di addestramento.
	
	La Figura~\ref{fig:dataset_distribution} illustra la distribuzione statistica dell'intero dataset e la proporzione rappresentata dal \textit{test set} di 85 liquidi, evidenziando la distribuzione per numero di componenti, i valori di viscosità esplorati e la presenza dei singoli costituenti chimici.
	
	\begin{figure}[h!]
		\centering
		% Sotto-figura (a)
		\begin{subfigure}[b]{0.495\textwidth}
			\centering \includegraphics[width=\textwidth]{numero_composti.png}
			\caption{Numero di composti nei liquidi.}
			\label{fig:dist_compounds}
		\end{subfigure}
		\hfill
		% Sotto-figura (b)
		\begin{subfigure}[b]{0.49\textwidth}
			\centering \includegraphics[width=\textwidth]{valori_viscosità.png}
			\caption{Distribuzione dei valori di viscosità.}
			\label{fig:dist_viscosity}
		\end{subfigure}
		
		% Sotto-figura (c)
		\begin{subfigure}[b]{\textwidth}
			\centering \includegraphics[width=\textwidth]{istogramma_composti.png}
			\caption{Numero di punti dati per singolo costituente chimico.}
			\label{fig:dist_elements}
		\end{subfigure}
		
		\caption{Analisi descrittiva del dataset globale (blu) e del sottoinsieme riservato al \textit{test set} (arancione). I dati mantengono le suddivisioni definite da Cassar per consentire un confronto omogeneo tra modelli diversi.~\cite{cassarViscNetNeuralNetwork2021}}
		\label{fig:dataset_distribution}
	\end{figure}
	
	\section{Estrazione delle features}
	\label{sec:feature_extraction}
	Per superare i limiti predittivi legati all'utilizzo delle sole frazioni molari, Cassar ha effettuato un'approfondita ingegnerizzazione delle \textit{features} (\textit{feature extraction}), trasformando il dominio composizionale in un dominio di proprietà chimico-fisiche~\cite{cassarViscNetNeuralNetwork2021}.
	
	Matematicamente, definendo $\boldsymbol{C} = [x_1, x_2, \dots, x_n]$ come il vettore delle frazioni molari atomiche degli elementi presenti nel liquido e $\boldsymbol{S} = [s_1, s_2, \dots, s_n]$ come il vettore di una specifica proprietà chimica per ciascun elemento (es. raggio atomico, elettronegatività, volume), si calcolano i vettori delle proprietà pesate ($\boldsymbol{W}$) e assolute ($\boldsymbol{A}$) attraverso il prodotto di Hadamard (moltiplicazione elemento per elemento):
	\begin{equation}
		\boldsymbol{W} = \boldsymbol{C} \circ \boldsymbol{S}
	\end{equation}
	e
	\begin{equation}
		\boldsymbol{A} = \lceil \boldsymbol{C} \rceil \circ \boldsymbol{S},
	\end{equation}
	dove l'operazione di \textit{ceiling} $\lceil \boldsymbol{C} \rceil$ agisce singolarmente su ogni componente di $\boldsymbol{C}$ e valuta esclusivamente la presenza o assenza dell'elemento, ignorandone la quantità relativa. A questi vettori vengono successivamente applicate funzioni di aggregazione statistiche (somma, media, deviazione standard, minimo e massimo) per estrarre i singoli descrittori scalari del liquido. Anche per questo lavoro di tesi si è applicata la stessa procedura sul medesimo catalogo di proprietà atomiche utlizzato da Cassar~\cite{cassarViscNetNeuralNetwork2021}, generando uno spazio ad alta dimensionalità composto da 560 \textit{features} iniziali.
	
	\section{La pipeline di machine learning}
	\label{sec:pipeline_myega}
	La struttura di base del modello predittivo adotta l'approccio \textit{gray-box} usato da Cassar per ViscNet, che è stato descritto alla fine della Sez.~\ref{sec:modelli_datadriven}. E possibile visualizzarne una rappresentazione schematica nella Fig.~\ref{fig:flowchart_viscnet}.
	Il distacco metodologico principale di questa tesi rispetto al lavoro di riferimento risiede nella scelta dell'architettura del modello. Laddove la pipeline originale impiega una rete neurale (ViscNet) a due \textit{layer} nascosti, il presente lavoro sostituisce tale architettura con modelli matematici strettamente lineari (Ridge, Elastic Net, VIF+OLS, PCR e SPCR).
	I modelli sono stati quindi allenati per prevedere i tre parametri fondamentali dell'equazione di viscosità MYEGA. Una volta ottenuti tali parametri, la viscosità $\eta$ a qualsiasi temperatura assoluta $T$ viene calcolata analiticamente tramite la seguente formulazione~\cite{mauroViscosityGlassformingLiquids2009}
	%
	\begin{equation}
		\begin{aligned}
			&\log_{10} \eta(T) \\
			&= \log_{10} \eta_\infty + (12 - \log_{10} \eta_\infty) \frac{T_g}{T} \exp \left[ \left( \frac{m}{12 - \log_{10} \eta_\infty} - 1 \right) \left( \frac{T_g}{T} - 1 \right) \right],
		\end{aligned}
		\label{eq:myega} 
	\end{equation}
	dove $T_g$ è la temperatura di transizione vetrosa, $m$ l'indice di fragilità e $\log_{10}\eta_\infty$ la viscosità asintotica in scala logaritmica (per maggiori dettagli vedere la Sez.~\ref{sec:liquidi_sottoraffreddati}).
	
	Lo scopo è di testare se, a fronte della medesima base dati e della medesima equazione MYEGA, uno spazio vettoriale lineare sia in grado di eguagliare la potenza predittiva del \textit{deep learning} garantendo al contempo una lettura fisica e trasparente dei pesi assegnati.
	
\end{document}